\chapter{Money and Inflation}\label{ch:b1:money-and-inflation}

In 1960 the United States' nominal GDP was $1.782$ times its stock of broad
money; in 1980 it was $1.855$ times, in 1990 $1.85$. By 2020 the ratio had
fallen to $1.182$. % figdata: velocity_us
The quantity equation $MV = PY$ held in every one of those years, because it
defines $V$. What failed is the idea that $V$ stays put, and with it the
simplest rule of monetary policy. This chapter states what money demand
predicts, scores the claims made for money growth against the record, and
follows money to its extreme, the \omterm{def:b1:money-and-inflation:hyperinflation}{hyperinflation}.

\begin{openingchart}{Velocity of broad money (nominal GDP divided by broad
money), 1950--2020.}{Jord\`a--Schularick--Taylor Macrohistory Database,
release 6}
\begin{tikzpicture}
\begin{axis}[omtime, xlabel={Year}, ylabel={Velocity}, xmin=1950, xmax=2020, ymin=0, ymax=3.6,
  width=0.72\linewidth, height=0.4\linewidth, legend style={font=\footnotesize, at={(0.5,-0.24)}, anchor=north,
  legend columns=4}]
  \addplot+[thick, no marks] table[x=year, y=usa, col sep=comma] {figdata/bachelor-1/23-money-and-inflation/out/velocity.csv};
  \addplot+[thick, no marks] table[x=year, y=gbr, col sep=comma] {figdata/bachelor-1/23-money-and-inflation/out/velocity.csv};
  \addplot+[thick, no marks] table[x=year, y=deu, col sep=comma] {figdata/bachelor-1/23-money-and-inflation/out/velocity.csv};
  \addplot+[thick, no marks] table[x=year, y=jpn, col sep=comma] {figdata/bachelor-1/23-money-and-inflation/out/velocity.csv};
  \legend{United States, United Kingdom, Germany, Japan}
\end{axis}
\end{tikzpicture}
\end{openingchart}

\omterm{def:b1:money-and-inflation:velocity}{Velocity} fell almost everywhere, and in Japan most: households and firms
hold more money per unit of income than they did. Only the American series
was roughly flat for decades, which is why it served as the textbook
example.

\begin{recall}
Money, its functions and its creation by banks were met at school, as were
the quantity equation, real and nominal interest rates and the Fisher
effect. \omterm{def:b1:money-and-inflation:velocity}{Velocity}, \omterm{def:b1:money-and-inflation:seigniorage}{seigniorage}, the \omterm{def:b1:money-and-inflation:seigniorage}{inflation tax} and \omterm{def:b1:money-and-inflation:hyperinflation}{hyperinflation} were
met there in words; this chapter defines them.
The \omterm{def:b1:long-run-model:neutrality}{neutrality of money} in the \omterm{def:b1:long-run-model:long-run}{long-run model} is
\cref{def:b1:long-run-model:neutrality}.
\end{recall}

\section{Money demand}\label{sec:b1:money-and-inflation:demand}

Money pays no interest, or less than bonds; holding it has an opportunity
cost, the nominal interest rate $i$. It is held because it buys things: the
more transactions, the more money.

\begin{definition}[Money demand function]\label{def:b1:money-and-inflation:money-demand}
\index{money demand function}\index{real money balances}
The \emph{money \omterm{def:b1:consumer-choice:demand-function}{demand function}} gives the real money balances people wish
to hold as a function of the nominal interest rate and real income,
$M/P = L(i, Y)$, decreasing in $i$ and increasing in $Y$.
\end{definition}

\begin{definition}[Velocity of money]\label{def:b1:money-and-inflation:velocity}
\index{velocity of money}
The \emph{velocity} of money is nominal output divided by the money stock,
$V = PY/M$: the number of times a unit of money is spent on final output in
a year. The quantity equation $MV = PY$ holds by this definition.
\end{definition}

\begin{notation}[Symbols of this chapter]\label{not:b1:money-and-inflation:symbols}
In this chapter $\beta$ is the semi-elasticity of money demand with respect to
the interest rate and $\alpha$ that of real balances with respect to expected
inflation (Cagan's money demand). Elsewhere in the series $\alpha$ is capital's
share and $\beta$ a discount factor or a regression coefficient.
\end{notation}

A common form is log-linear, $\ln(M/P) = a + \eta \ln Y - \beta i$, with an income
elasticity $\eta$ and a semi-elasticity $\beta$ with respect to the interest rate.
\omterm{def:b1:money-and-inflation:velocity}{Velocity} is then $V = PY/M = Y/L(i, Y)$; with $\eta = 1$, $\ln V = -a + \beta i$: \omterm{def:b1:money-and-inflation:velocity}{velocity} rises
with the interest rate, and is constant only if $i$ is.

\begin{omfigure}
\begin{tikzpicture}
\begin{axis}[omrecord, grid=none, xlabel={Real money balances $M/P$}, ylabel={Interest rate $i$ (\%)},
  xmin=0, xmax=10, ymin=0, ymax=10, width=0.5\linewidth, height=0.4\linewidth, samples=60]
  \addplot[demand, domain=1:10] {40/x - 3};
  \addplot[supply] coordinates {(5,0) (5,10)};
  \addplot[shifted] coordinates {(4,0) (4,10)};
  \addplot[guide] coordinates {(0,5) (5,5)};
  \addplot[guide] coordinates {(0,7) (4,7)};
  \node[font=\footnotesize, anchor=west] at (axis cs:8,1.2) {$L(i, Y)$};
  \node[font=\footnotesize, anchor=south west] at (axis cs:5,9) {$M/P$};
\end{axis}
\end{tikzpicture}

{\small The money market: real money demand (blue) and the real money supply
$M/P$ (red). A smaller real money supply (dashed) raises the interest rate that
makes people content to hold it. In the long run the interest rate is set
elsewhere and $P$ adjusts to bring $M/P$ to $L(i, Y)$.}
\end{omfigure}

\begin{measure}[Broad and narrow money]
Narrow money is currency and deposits that can be spent at once; broad
money adds savings and time deposits and similar claims. The Macrohistory
Database gives both, in local currency, for 18 advanced economies annually
since 1870 (release 6). % ledger: b1:jst-money-docs
The \omterm{def:b1:money-and-inflation:velocity}{velocity} of the opening chart uses broad money; a country where people
move savings into deposits sees its broad-money \omterm{def:b1:money-and-inflation:velocity}{velocity} fall although
nothing has changed in how they pay.
\end{measure}

\section{The quantity theory}\label{sec:b1:money-and-inflation:quantity}

The quantity equation $MV = PY$ becomes a theory when something is said
about $V$.

\begin{definition}[Quantity theory of money]\label{def:b1:money-and-inflation:quantity-theory}
\index{quantity theory of money}
The \emph{quantity theory of money} is the quantity equation together with
the hypothesis that \omterm{def:b1:money-and-inflation:velocity}{velocity} is stable (or moves predictably) and that real
output is set by real factors; then the price level is proportional to the
money stock, and in growth rates $\pi = g_M + g_V - g_Y$, with $g_V$ near zero.
\end{definition}

\begin{omchart}{Average inflation against average growth of broad money, 18
advanced economies, decades 1870--2019 (per cent a year), with the fitted
line.}{Jord\`a--Schularick--Taylor Macrohistory Database, release 6}
\begin{tikzpicture}
\begin{axis}[omscatter, xlabel={Money growth (\% a year)}, ylabel={Inflation (\% a year)}, xmin=-5, xmax=40,
  ymin=-10, ymax=55, width=0.6\linewidth, height=0.45\linewidth]
  \addplot[only marks, mark=*, mark size=1.3pt, omDef, fill opacity=0.6] table[x=gm, y=pi, col sep=comma] {figdata/bachelor-1/23-money-and-inflation/out/qt_decades.csv};
  \addplot[omThm, thick, sharp plot, mark=none] table[x=gm, y=fit, col sep=comma] {figdata/bachelor-1/23-money-and-inflation/out/qt_fit.csv};
\end{axis}
\end{tikzpicture}
\end{omchart}

Over 228 country-decades the correlation between money growth and
inflation is $0.848$, and inflation rises by $0.931$ points for each point of
money growth; the intercept, $-3.24$, is roughly minus the growth of output.
% figdata: qt_stats
The slope leans on the high-inflation decades. Without the two wartime
decades at the top right, Italy and Japan in the 1940s, it is $0.758$; over
the 206 decades with money growth below $15\%$ a year it is $0.605$, and over
the 123 decades since 1950, $0.559$. % figdata: qt_stats
Over 110 countries and 30 years the same correlation is high for every
definition of money (McCandless and Weber 1995). % ledger: b1:mccandless1995
Over a decade, countries with fast money growth have high inflation. Within
a decade, and at low inflation, the link is loose: \omterm{def:b1:money-and-inflation:velocity}{velocity} moves.

\begin{model}{The quantity theory}
\assumes \omterm{def:b1:money-and-inflation:velocity}{Velocity} stable or predictable; output set by real factors.
\predicts Inflation equal to money growth minus output growth in the long
run.
\tested Across countries and decades, money growth and inflation are
highly correlated (the chart; McCandless and Weber
1995). % ledger: b1:mccandless1995
The slope is near one only when high-inflation decades are included:
$0.931$ with them, $0.6$ or less at moderate money growth. % figdata: qt_stats
\breaks At low inflation \omterm{def:b1:money-and-inflation:velocity}{velocity} moves, as the opening chart shows, and
policy aimed at money growth was given up in the 1980s.
\end{model}

\section{The money-growth propositions scored}\label{sec:b1:money-and-inflation:scored}

Three claims were long made for money growth. Each can be scored.

\emph{(i) Over the long run, money growth and inflation move one for one.}
Holds for high inflation, not at low inflation: the decade slope is $0.758$
to $0.931$, driven by the high-inflation decades, and well below one
($0.605$) where money grew by less than $15\%$ a year. % figdata: qt_stats

\emph{(ii) Money affects output and prices with long and variable delays.}
Holds as a description: the effects of monetary policy on output come
first and on prices later, as Part~VI and the second university year
measure.

\emph{(iii) A central bank can stabilise inflation by keeping money growth
on a target.} Failed when \omterm{def:b1:money-and-inflation:velocity}{velocity} moved.

\begin{record}[Velocity in the 1980s]
On 6 October 1979 the Federal Reserve shifted its focus from the
overnight interest rate to the volume of bank reserves, to gain firmer
control over the growth of the money supply; its earlier procedure had
become less reliable. The federal funds rate reached a record $20\%$ in late
1980. In July 1982 the chairman told Congress that the Federal Reserve was
backing off its previous tight targets. % ledger: b1:fed-targets
In the American data, broad-money \omterm{def:b1:money-and-inflation:velocity}{velocity} was $1.855$ in 1980 and $1.85$ in
1990, flat over the decade, but it rose to $2.14$ by 2000 and then fell to
$1.182$ in 2020. % figdata: velocity_us
A target for money growth is a target for inflation only if \omterm{def:b1:money-and-inflation:velocity}{velocity}
cooperates.
\end{record}

\section{The inflation tax and hyperinflation}\label{sec:b1:money-and-inflation:tax}

\begin{definition}[Seigniorage and the inflation tax]\label{def:b1:money-and-inflation:seigniorage}
\index{seigniorage}\index{inflation tax}
\emph{Seigniorage} is the real revenue a government obtains by issuing
money, $\dot M/P$. The \emph{inflation tax} is the loss of real value that
inflation inflicts on money holders, $\pi M/P$: a tax at the rate $\pi$ on the
base $M/P$. With no \omterm{prop:b1:models-and-data:real-growth}{real growth}, in a \omterm{def:b1:solow-model:steady-state}{steady state} where money grows at the
rate of inflation, the two are equal.
\end{definition}

In a \omterm{def:b1:solow-model:steady-state}{steady state} with real balances $m = M/(PY)$ (a share of output) and
money growing at the rate $\pi$ (no \omterm{prop:b1:models-and-data:real-growth}{real growth}), \omterm{def:b1:money-and-inflation:seigniorage}{seigniorage} as a share of
output is $\pi m$. But the base shrinks as the rate rises.

\begin{definition}[Hyperinflation]\label{def:b1:money-and-inflation:hyperinflation}
\index{hyperinflation}
A \emph{hyperinflation} is an episode in which prices rise by more than
$50\%$ a month (Cagan's threshold, about $13{,}000\%$ a
year). % ledger: b1:cagan-def
\end{definition}

\begin{model}{Money demand in a hyperinflation}
\assumes Real balances fall exponentially with expected inflation,
$m = k e^{-\alpha\pi^e}$ (Cagan's money demand), and in a \omterm{def:b1:solow-model:steady-state}{steady state} $\pi^e = \pi$.
\predicts \omterm{def:b1:money-and-inflation:seigniorage}{Seigniorage} $\pi k e^{-\alpha\pi}$, largest at $\pi = 1/\alpha$; a government that needs more
cannot get it, and inflation explodes.
\tested Consistent with it: real balances fell as inflation rose in Bolivia
in 1984--85 (the chart below). A test would estimate $\alpha$ and compare the
inflation rates governments chose with $1/\alpha$.
\breaks When inflation accelerates, the \omterm{def:b1:solow-model:steady-state}{steady state} $\pi^e = \pi$ fails; and
when holders can protect their money, as in Peru, where broad money barely
fell through a \omterm{def:b1:money-and-inflation:hyperinflation}{hyperinflation} (the record below).
\end{model}

\begin{proposition}[Revenue-maximising inflation]\label{prop:b1:money-and-inflation:laffer}
With $m = ke^{-\alpha\pi}$, steady-state \omterm{def:b1:money-and-inflation:seigniorage}{seigniorage} $S(\pi) = \pi k e^{-\alpha\pi}$ is maximal at $\pi^* = 1/\alpha$, where
it equals $k/(\alpha e)$ of output. It does not depend on the level of output.
\end{proposition}

\begin{proof}
$S'(\pi) = ke^{-\alpha\pi}(1 - \alpha\pi)$, positive below $1/\alpha$ and negative above. At $\pi^* = 1/\alpha$,
$S = (1/\alpha)ke^{-1}$. Output enters $m$ only as the scale of $M/P$, and $S$ is a share of it.
\end{proof}

\begin{omfigure}
\begin{tikzpicture}
\begin{axis}[omrecord, grid=none, xlabel={Inflation $\pi$ (continuous rate)}, ylabel={Seigniorage (\% of GDP)},
  xmin=0, xmax=2, ymin=0, ymax=4.5, width=0.56\linewidth, height=0.38\linewidth, samples=100]
  \addplot[demand, domain=0:2] {100*x*0.25*exp(-2.5*x)};
  \addplot[guide, domain=0:2] {2};
  \addplot[guide] coordinates {(0.4,0) (0.4,3.679)};
  \node[font=\footnotesize, anchor=south west] at (axis cs:0.42,3.68) {$3.68\%$ at $\pi = 0.4$};
  \node[font=\footnotesize, anchor=north west] at (axis cs:0.11,1.95) {$0.104$};
  \node[font=\footnotesize, anchor=south] at (axis cs:1.0171,2.05) {$1.017$};
\end{axis}
\end{tikzpicture}

{\small The Laffer curve of the \omterm{def:b1:money-and-inflation:seigniorage}{inflation tax} in the weekend problem,
$m = 0.25e^{-2.5\pi}$. A deficit of $2\%$ of GDP (dashed) can be financed at a low
inflation rate, $0.104$, or a high one, $1.017$; no rate finances more than $3.68\%$.}
\end{omfigure}

\begin{omchart}{Broad money (\% of GDP) against consumer-price inflation (\% a
year, log scale), Bolivia and Peru, 1980--1996.}{World Bank, World Development
Indicators (FM.LBL.BMNY.GD.ZS, FP.CPI.TOTL.ZG)}
\begin{tikzpicture}
\begin{axis}[omscatter, xmode=log, log ticks with fixed point, xmin=5, xmax=20000, ymin=0, ymax=55,
  xlabel={Inflation (\% a year, log scale)}, ylabel={Broad money (\% of GDP)}, width=0.62\linewidth,
  height=0.42\linewidth, legend style={font=\footnotesize, at={(0.97,0.97)}, anchor=north east}]
  \addplot[omDef, thin, sharp plot, mark=*, mark size=1.4pt] table[x=infl, y=money, col sep=comma] {figdata/bachelor-1/23-money-and-inflation/out/hyper_bol.csv};
  \addplot[omThm, thin, sharp plot, mark=square*, mark size=1.3pt] table[x=infl, y=money, col sep=comma] {figdata/bachelor-1/23-money-and-inflation/out/hyper_per.csv};
  \legend{Bolivia, Peru}
\end{axis}
\end{tikzpicture}
\end{omchart}

\begin{record}[Bolivia 1985]
Bolivia's inflation was $123.5\%$ in 1982, $1{,}281.3\%$ in 1984 and $11{,}749.6\%$ in
1985; broad money fell from $27.1\%$ of GDP in 1982 to $12.1\%$ in 1985 and
$10.7\%$ in 1986. After the stabilisation inflation was $14.6\%$ in 1987, and
broad money climbed back to $46\%$ of GDP by 1994. % figdata: hyper_bol
People held money again once it kept its value. Peru's broad money barely
moved through its own \omterm{def:b1:money-and-inflation:hyperinflation}{hyperinflation}, $24.1\%$ of GDP in 1990 with inflation
of $7{,}481.7\%$ --- a reminder that a broad aggregate can include deposits
that are protected from inflation. % figdata: hyper_per
\end{record}

\begin{omfigure}
\begin{tikzpicture}
\begin{axis}[omtime, grid=none, xlabel={Time (years)}, ylabel={Logarithm}, xmin=0, xmax=10, ymin=0, ymax=0.6,
  width=0.52\linewidth, height=0.36\linewidth, legend style={font=\footnotesize, at={(0.03,0.97)}, anchor=north west}]
  \addplot[demand, domain=0:10, samples=11] {x < 4 ? 0.05 + 0.02*x : 0.13 + 0.06*(x - 4)};
  \addplot[supply, domain=0:3.999, samples=10] {0.02*x};
  \addplot[supply, forget plot, domain=4:10, samples=10] {0.18 + 0.06*(x - 4)};
  \addplot[guide] coordinates {(4,0) (4,0.6)};
  \legend{$\ln M$, $\ln P$}
\end{axis}
\end{tikzpicture}

{\small Money growth rises from $2\%$ to $6\%$ a year at year 4 (dashed line); $\ln M$
is drawn with an arbitrary vertical offset. With
$\beta = 2.5$, the expected rise of $4$ points in inflation lowers money demand by $0.1$ in
logarithm (about $10\%$):
the price level jumps by $0.1$ in logarithm at once, then grows at the new
rate. Prices run ahead of money.}
\end{omfigure}

\section{The Fisher effect and the costs of inflation}\label{sec:b1:money-and-inflation:fisher}

The nominal interest rate is $i = r + \pi^e$. If the real rate is set by real
factors (\cref{ch:b1:long-run-model}), expected inflation passes one for one
into nominal rates: the Fisher effect.

Over decades since 1960 that is roughly what the data show: across the
advanced economies, short-term interest rates rose by $0.871$ points for
each point of decade inflation, with a correlation of $0.761$ over 105
country-decades. Over all decades since 1870 the slope is only $0.223$.
% figdata: qt_stats
Under the gold standard, inflation was not expected to persist --- prices
rose and fell --- and the interest rate did not follow it.

Inflation has costs even when expected. People hold less money and spend
time and resources economising on it (shoe-leather costs); prices must be
reprinted (menu costs); a tax system written in nominal terms taxes
nominal gains and nominal interest (tax distortions); and inflation that
is high is also variable, which makes long contracts risky.

\section*{Exercises}\label{sec:b1:money-and-inflation:exercises}

\begin{exercise}[$\star$]\label{exo:b1:money-and-inflation:qt}
Money grows by $8\%$ a year, output by $3\%$, and \omterm{def:b1:money-and-inflation:velocity}{velocity} is constant. Find
inflation.
\end{exercise}

\begin{exercise}[$\star$]\label{exo:b1:money-and-inflation:read}
American broad-money \omterm{def:b1:money-and-inflation:velocity}{velocity} was $1.782$ in 1960, $2.14$ in 2000 and $1.182$ in
2020. % figdata: velocity_us
With \omterm{def:b1:money-and-inflation:velocity}{velocity} at its 1960 level, by what factor would nominal GDP in 2020
have exceeded its actual value, for the same money stock?
\end{exercise}

\begin{exercise}[$\star$]\label{exo:b1:money-and-inflation:seign}
Inflation is $20\%$ a year and real balances are $10\%$ of GDP. Compute the
\omterm{def:b1:money-and-inflation:seigniorage}{inflation tax}.
\end{exercise}

\begin{exercise}[$\star$]\label{exo:b1:money-and-inflation:fisher-ex}
Expected inflation rises by 3 points and the real rate is unchanged. Predict
the nominal rate.
\end{exercise}

\begin{exercise}[$\star\star$]\label{exo:b1:money-and-inflation:elast}
With $\ln(M/P) = a + \ln Y - 4i$, find the elasticity of \omterm{def:b1:money-and-inflation:velocity}{velocity} with respect to $i$ at
$i = 0.05$.
\end{exercise}

\begin{exercise}[$\star\star$]\label{exo:b1:money-and-inflation:max}
With $M/P = Ye^{-2\pi}$, find the revenue-maximising inflation rate and the
maximum \omterm{def:b1:money-and-inflation:seigniorage}{seigniorage} as a share of $Y$.
\end{exercise}

\begin{exercise}[$\star\star$]\label{exo:b1:money-and-inflation:bolivia}
From the Bolivia record, broad money was $27.1\%$ of GDP at inflation of
$123.5\%$ (1982) and $12.1\%$ at $11{,}749.6\%$ (1985). % figdata: hyper_bol
Compute the ratio of real balances. What does it say about the base of the
\omterm{def:b1:money-and-inflation:seigniorage}{inflation tax}?
\end{exercise}

\begin{exercise}[$\star\star$]\label{exo:b1:money-and-inflation:velfall}
Money grows by $6\%$ a year, output by $2\%$, and \omterm{def:b1:money-and-inflation:velocity}{velocity} falls by $1\%$ a year.
Predict inflation.
\end{exercise}

\begin{exercise}[$\star\star$]\label{exo:b1:money-and-inflation:fisher-slope}
The slope of decade interest rates on decade inflation is $0.223$ over all
decades since 1870 and $0.871$ since 1960. % figdata: qt_stats
What does each say about the Fisher effect?
\end{exercise}

\begin{exercise}[$\star\star\star$]\label{exo:b1:money-and-inflation:derive}
Prove \cref{prop:b1:money-and-inflation:laffer}, and show that with \omterm{prop:b1:models-and-data:real-growth}{real growth}
at the rate $g$, \omterm{def:b1:money-and-inflation:seigniorage}{seigniorage} is $(\pi + g)m$; find the inflation rate that maximises it.
\end{exercise}

\begin{exercise}[$\star\star\star$]\label{exo:b1:money-and-inflation:rule}
A central bank keeps money growth at $g_M$; output grows at $g_Y$; \omterm{def:b1:money-and-inflation:velocity}{velocity} growth
is $\bar g_V + \varepsilon_t$ with $\varepsilon_t$ independent, of mean zero and variance $\sigma^2$. Find the mean and
variance of inflation, and compare with a bank that observes $\varepsilon_t$ and offsets
it.
\end{exercise}

\begin{exercise}[$\star\star\star$]\label{exo:b1:money-and-inflation:refute}
What evidence would refute ``inflation is always caused by money growth'' at
low inflation, and why does the cross-country, cross-decade link survive
it?
\end{exercise}

\section{Problem: Kestria's Printing Press}\label{sec:b1:money-and-inflation:problem}

\begin{problem}[{Weekend problem --- the inflation tax}]\label{pb:b1:money-and-inflation:kestria}
Kestria (invented) has no \omterm{prop:b1:models-and-data:real-growth}{real growth}. Its \omterm{def:b1:money-and-inflation:money-demand}{real money balances}, as a share
of GDP, are $m = 0.25e^{-2.5\pi}$, where $\pi$ is the continuously compounded inflation
rate, expected correctly. In a \omterm{def:b1:solow-model:steady-state}{steady state} money grows at the rate $\pi$.

\medskip\noindent\textbf{Part I --- Money demand and velocity.}
\begin{enumerate}
  \item Compute $m$ and \omterm{def:b1:money-and-inflation:velocity}{velocity} at $\pi = 0.05$.
  \item Compute them at $\pi = 0.10$.
  \item What is the semi-elasticity of money demand, and the elasticity at $\pi = 0.05$?
  \item Compute \omterm{def:b1:money-and-inflation:seigniorage}{seigniorage} at $\pi = 0.05$.
  \item Why does \omterm{def:b1:money-and-inflation:velocity}{velocity} rise with inflation?
\end{enumerate}

\medskip\noindent\textbf{Part II --- A deficit paid by printing.} The government
must finance a deficit of $2\%$ of GDP by money creation.
\begin{enumerate}[resume]
  \item Write the condition for the steady-state inflation rate.
  \item Find the two rates that satisfy it.
  \item Compute real balances at each.
  \item Which is the government likely to reach if it starts at low inflation?
  \item What happens if people expect the high rate?
\end{enumerate}

\medskip\noindent\textbf{Part III --- The Laffer curve.}
\begin{enumerate}[resume]
  \item Find the revenue-maximising inflation rate.
  \item Express it as an annual percentage increase in prices.
  \item Compute the maximum \omterm{def:b1:money-and-inflation:seigniorage}{seigniorage}.
  \item Compute \omterm{def:b1:money-and-inflation:seigniorage}{seigniorage} at $\pi = 0.6$ and at $\pi = 1.5$.
  \item The deficit rises to $5\%$ of GDP. What happens?
  \item Why does the maximum not depend on the size of Kestria's GDP?
\end{enumerate}

\medskip\noindent\textbf{Part IV --- A stabilisation.} Kestria is at $\pi = 1.5$. A
fiscal reform ends the deficit and a new currency is introduced; inflation
falls at once to $\pi = 0.02$.
\begin{enumerate}[resume]
  \item Compute real balances before, and the annual price increase at $\pi = 1.5$.
  \item Compute real balances after.
  \item By what factor do real balances rise?
  \item How much new money, as a share of GDP, can the central bank issue without raising prices?
  \item Why must the fiscal reform come first?
  \item Compare with Bolivia after 1985.
  \item What would make people doubt the stabilisation, and what would that do to real balances?
  \item State the result: what is the largest deficit Kestria can finance by printing?
\end{enumerate}
\end{problem}
